\chapter{Tabla de dominio, 1909--1945}
\label{cap:dominio}

Este apéndice reúne los actos sobre el dominio de la tierra del departamento que aparecieron
en el barrido del Boletín Oficial entre 1909 y diciembre de 1945: deslindes, mensuras,
remates judiciales, reposiciones de título y juicios de posesión treintañal. Son
\textbf{sesenta y tres}, y están ordenados por fecha.

No es un registro de propiedad ni pretende serlo. \textbf{Es la lista de las veces en que el
dominio de un pedazo de este territorio tuvo que pasar por un aviso público}, y por eso dice
menos sobre quién era dueño de qué que sobre cómo se acreditaba, se partía y se perdía la
tierra en esos treinta y seis años.

\begin{cautela}
\textbf{Todo lo que está aquí sale de edictos y avisos, no de folios registrales.} Un remate
publicado no prueba que la subasta se realizara ni quién compró; un deslinde iniciado no prueba
que se terminara; una posesión treintañal deducida no prueba que se concediera. Los nombres son
los que el aviso imprime, con las homonimias que eso arrastra. Y la columna de superficie está
casi vacía porque \textbf{la mayoría de los avisos no la dan}: describen por linderos.

\textbf{Cada fila dice dónde se lee, cuando se pudo.} Las \textbf{treinta y dos filas de 1909 a 1945} que salen de tramos leídos edición por edición llevan al final de la columna de partes el \textbf{número de Boletín y la hoja} en que el aviso salió, de modo que cualquiera puede ir a verlo. Las \textbf{treinta y una restantes, de 1924 a 1945, todavía no}: se armaron a partir de un barrido cuyo registro no consigna la edición fila por fila, y localizarlas exige rehacer el año. \textbf{Completarlas es trabajo pendiente de este libro, y está declarado en el apéndice \ref{cap:pedidos}.}

\textbf{Y el barrido tiene huecos.} Faltan treinta y un números del tramo 1923--1945 (apéndice
\ref{cap:fuentes}): lo que no figura aquí puede estar en ellos. Entre 1922 y 1923 la lectura fue parcial. En 1920, leído edición por edición al cierre, \textbf{tampoco hay un solo acto de dominio}: los diez actos del año son de policía, de registro civil y de elecciones. \textbf{En 1921, en cambio, vuelven}: dos deslindes, los primeros del departamento desde 1917, y son las dos filas de ese año. En el \textbf{primer semestre de 1922}, leído edición por edición al cierre, \textbf{no hay acto de dominio} salvo el aviso de remate del Potrero del Gallinato, que es la fila de ese año y de la que el archivo \textbf{no dice el resultado}: la subasta se fijó para seis días después del final del tramo leído. \textbf{En 1923, leído entero, vuelve un acto de dominio y es el primero desde 1921}: el remate de la \textbf{finca Calderilla} en el juicio sucesorio de Silvano Murúa\index{Murúa, Silvano Ignacio}, con base de \$7.000 y subasta fijada para el 11 de diciembre, trece días después del cierre del año; tampoco de éste dice el archivo si se vendió. En 1919, leído edición por edición al cierre, \textbf{tampoco hay un solo acto de dominio}: los nueve actos del año son de policía, de gobierno local, de registro civil y de presupuesto. En 1918, leído edición por edición al cierre, \textbf{tampoco hay un solo acto de dominio del departamento}: los veintisiete actos del año son de policía, de agua, de escuela, de minas y de elecciones.  En 1915, leído edición por edición al cierre, \textbf{no hay un solo acto de dominio del departamento}; en 1916, leído del mismo modo, hay \textbf{tres}, y son las tres filas de ese año. Las tres filas de 1914 salen del barrido edición por edición del año entero, hecho en septiembre de 2026, al que le faltan las siete ediciones consecutivas de junio y una más, la Nº 470 (apéndice \ref{cap:fuentes}).

Para convertir esta tabla en una cadena de dominio hacen falta los folios que el apéndice
\ref{cap:pedidos} reclama. Lo que la tabla sí permite es \textbf{saber qué pedir}, finca por
finca y con fecha.
\end{cautela}

{\footnotesize
\begin{longtable}{@{}p{0.8cm}p{2.4cm}p{1.6cm}p{2.1cm}p{2.9cm}p{2.1cm}@{}}
\toprule
\textbf{Año} & \textbf{Finca o inmueble} & \textbf{Paraje} & \textbf{Acto} & \textbf{Partes y linderos} & \textbf{Superficie o base} \\
\midrule
\endhead
1909 & San Alejo y Santa Rufina & San Alejo & \textbf{Deslinde general} & División de condominio: Benítez y Aguilar contra la familia Fernández\index{Fernández, Ángel}; perito Walter Hesling. Linderos: al sud Garzón y Pinto, herederos de Desiderio Ruiz y Silverio Bejarano; al naciente Daniel Linares\index{Linares, Daniel}, Abelardo Figueroa\index{Figueroa, Abelardo} y Nolasco Cornejo; al norte Francisco Sasaya, Domingo Royo y la finca Cerro Negro (Jujuy); al poniente el Cerro Nevado. Continúa con incidentes en 1910 \textit{(B.O.\ Nº 59, h.\ 4)}& --- \\
\addlinespace
1909 & Angosto de Arrieta & Angosto de Arrieta & Remate de \textbf{una décima parte indivisa} & Linderos: al este Julio Mercado, Augusto Rejis\index{Regis, Augusto} y Susana Aguilera; al sud la estancia Quitilipe; al oeste el camino nacional y Nolasco J.\ Cornejo; al norte Cornejo y Mercado \textit{(B.O.\ Nº 92, h.\ 4)}& Base \$300; finca tasada en \$3.000 \\
\addlinespace
1909 & Sucesión Catacata & Potrero de Castillo & Remate de haciendas, \textbf{dos veces} (febrero y diciembre) & Entrega en el paraje; martillero Manuel Nuñez de la Rosa. En febrero, 353 cabezas de vacuno y 124 de ovino; el detalle no cierra con los totales impresos \textit{(B.O.\ Nº 33, h.\ 4, y Nº 113, h.\ 4)}& Sin base, al contado \\
\addlinespace
1910 & \textbf{Los Porongos} & Los Porongos & \textbf{Deslinde, mensura y amojonamiento} & Pedido por Francisco Alemán, con poder de Mónica Diez de Costas. Linderos: al poniente, lomas que lindan con «tierras de la Caldera y Sauces»; al norte, tierras de Perico (herederos de Domingo Iriarte); al naciente, la loma «Las Víboras» y la Cañada Seca, hasta Los Piñones; al sud, «Tabacal» o «Pelada Muerta», que la divide de La Despensa \textit{(B.O.\ Nº 140, h.\ 4, repetido en Nº 190, h.\ 4)}& --- \\
\addlinespace
1910 & Loma Bola & --- & \textbf{Deslinde, mensura y amojonamiento} (nuevo día) & Pedido por Francisco Alemán por Ángel Costaz; agrimensor Rafael M.\ Zuvíria. Linderos: al norte, sucesores del doctor Julio Arias; al sud, herederos de Segundo D.\ Bedoya; al este, la estancia Los Porongos; al oeste, Francisco Sosaya \textit{(B.O.\ Nº 189, h.\ 4)}& --- \\
\addlinespace
1910 & Sucesión Catacata & Potrero de Castillo & Remate del \textbf{derecho de marca} y de los semovientes & Martillero M.\ Nuñez de la Rosa; tercer remate de la misma sucesión \textit{(B.O.\ Nº 174, h.\ 4)}& Sin base \\
\addlinespace
1911 & \textbf{San Alejo y Santa Rufina} & San Alejo & \textbf{Remate} de la finca entera, en la división de condominio & Herederos Fernández\index{Fernández, Ángel}; juez Vicente Arias, martillero Víctor M.\ Saravia. Mensura de Walter Hessling aprobada por el Departamento Topográfico; su honorario, \$5.000, se aprueba en junio. Linderos: norte, Jujuy, cabaña de Marcos Zorrilla, Uracatao o Pueblo Viejo de Domingo Royo y Gavino Sánchez; este, finca Angostura de Daniel Linares\index{Linares, Daniel}, Eusebio Herrera, Francisco Sajama, Arturo Escudero y Abelardo Figueroa\index{Figueroa, Abelardo}; sud, Linares, herederos de Desiderio Ruiz, Garzón y Pintos y Silverio Bejarano; oeste, José Bruzzo y Zorrilla \textit{(B.O.\ Nº 306, h.\ 4)}& \textbf{17.874 ha 27 a 79}; base \$180.000 \\
\addlinespace
1911 & \textbf{Acheral} & --- & Remate judicial (juez Alejandro Bassani; subasta del 18 de setiembre) & Martillero M.\ R.\ Alvarado \textit{(B.O.\ Nº 276, h.\ 4)}& Base \$4.000 \\
\addlinespace
1912 & \textbf{San Alejo y Santa Rufina} & San Alejo & \textbf{Segundo remate} (abril), tras no venderse en diciembre & Mismo juicio y martillero; el aviso la llama «la finca más importante del Depto.\ de la Caldera» y declara bosques, pasto para 8.000 vacunos y \textbf{una Escuela Nacional} en la finca. Garzón y Pintos figuran aquí al oeste y en 1911 al sud \textit{(B.O.\ Nº 328, h.\ 4)}& 17.874 ha; base \textbf{\$135.000} \\
\addlinespace
1912 & \textbf{Los Porongos} o \textbf{San Luis} & Los Porongos & Deslinde, mensura y amojonamiento & Pedido por Pablo Enrique Wernner y Enrique Leconte; peritos Hermán Pfister y Rodolfo Martín; juez Alejandro Bassani. Mismos linderos que el edicto de 1910, con «los Peñales» en lugar de Los Piñones \textit{(B.O.\ Nº 364, h.\ 4)}& --- \\
\addlinespace
1912 & \textbf{Santa Mónica} y \textbf{Quitilipe} & --- & Deslinde, mensura y amojonamiento & Pedido por José Cobach; perito Manuel Lapido. Santa Mónica linda al poniente con Yacones y San Alejo, al naciente con el camino nacional a Jujuy «cercado de piedras formando callejón» y al norte con el \textbf{río de la Caldera}; Quitilipe, al sud y poniente con el mismo río \textit{(B.O.\ Nº 372, h.\ 3)}& --- \\
\addlinespace
1913 & \textbf{Finca Lesser} & Lesser & Remate judicial (4 de junio, confitería del Jockey Club, capital) & Mart.\ Víctor M.\ Saravia. El aviso la vende como paisaje: «el punto más pintoresco de la provincia», casa de campo, quinta, «a una hora de viaje en carruaje de la ciudad». Linderos al sud, naciente y poniente con propiedad que fue de Narciso Figueroa, hoy de Arquati \textit{(B.O.\ Nº 399, h.\ 4)}& 1.600 cuadras cuadradas \\
\addlinespace
1913 & \textbf{Casa-quinta del pueblo} & Pueblo de La Caldera & Remate judicial de la sucesión Peralta (1º de julio, capital) & Mart.\ José María Leguizamón; juez Vicente Arias. Siete habitaciones sobre una cuadra cuadrada; al naciente la calle pública y por el poniente, sur y norte, propiedad de Jacinto Guerrero\index{Guerrero, Jacinto} \textit{(B.O.\ Nº 409, h.\ 3)}& Base \$1.500 \\
\addlinespace
1913 & \textbf{Terreno de Pablo Enrique Wenner} & --- & Deslinde, mensura y amojonamiento & Juez Alejandro Bassani ---«Baassani» \textit{[sic]} en el edicto---; peritos Pfister y Martín. Linderos: norte, Jacoba Molina; sur, Lucas Molina; naciente, la cumbre que divide con el Potrero de Valencia; poniente, el «Río de la Caladera» \textit{[sic]}. El auto manda publicarlo «por una vez» y el Boletín lo publica dos \textit{(B.O.\ Nº 443, h.\ 4, y Nº 449, h.\ 4)}& --- \\
\addlinespace
1914 & \textbf{San Alejo} y \textbf{Santa Rufina} & San Alejo & \textbf{Remate tres veces} (21 de febrero, 12 de mayo y 5 de octubre) & Mart.\ M.\ R.\ Alvarado, por orden del juez Vicente Arias; división de condominio de los herederos Fernández\index{Fernández, Ángel}. B.O.\ Nº 463, 476 y 509 \textit{(B.O.\ Nº 463, h.\ 3; Nº 476, h.\ 4; Nº 509, h.\ 3)}& \textbf{17.874 ha 27 a 19 m\textsuperscript{2}}; bases de \$5,50, \$4,25 y \$3,20 la hectárea \\
\addlinespace

1914 & \textbf{Campo Alegre} & --- & \textbf{Deslinde, mensura y amojonamiento} & \textbf{Francisco Sosaya}, con «títulos bastantes»; juez Francisco F.\ Sosa; agrim.\ Ettore Chiostri; secretario Nolasco Zapata. Linderos: norte, los Mercado; naciente, terrenos que fueron de Abelino Costas; sur, los mismos sucesores de Costas y de José Fernández; oeste, los sucesores Fernández. B.O.\ Nº 512 \textit{(B.O.\ Nº 512, h.\ 4)}& Sin superficie \\
\addlinespace

1914 & \textbf{Angostura} & --- & Remate judicial (juez Dr.\ Francisco F.\ Sosa; subasta del 23 de diciembre) & Mart.\ Ricardo López; linderos: norte, Francisco Sosaya; este, finca Loma Bola y propiedades de Julio Mercado y Lauredo Peralta; sur, Abelardo Figueroa; oeste, herederos de José Manuel Fernández. B.O.\ Nº 521 \textit{(B.O.\ Nº 521, h.\ 3)}& Base \$20.000, dos tercios de la tasación fiscal de \$30.000 \\
\addlinespace
1916 & \textbf{Loma Bola} & --- & \textbf{Deslinde, mensura y amojonamiento} & Pedido por Alberto Escudero; perito Herman Pfister\index{Pfister, Hermann}; juez Martín Gómez Rincón, secretario Nolasco Zapata. Linderos: norte, La Angostura de Daniel Linares\index{Linares, Daniel}; este, Los Porongos «\textit{que fue de Mónica Díaz de Cortes, hoy del conde Garnier de Turava}»; sud, Julio Mercado; oeste, el arroyo de la Angostura, Francisco Sosaya y Celedonia Reyes del Prado. Publicado dos veces \textit{(B.O.\ Nº 592, h.\ 4, y Nº 593, h.\ 3)}& --- \\
\addlinespace
1916 & \textbf{Los Porongos} hoy \textbf{San Luis} & Los Porongos & \textbf{Remate judicial} (juez Martín Gómez Rincón) & Ejecución de Enrique Werner contra \textbf{Humberto Graf von Garnier Turawa}; martillero Enrique Sylvester; subasta fijada para el 15 de junio en el Jockey Bar. Linderos: tierras de la Caldera y Sauces, herederos de Domingo Iriarte, estancia de los Peñones, «\textit{Tabacal}» o «\textit{Pelada Muerta}» y La Despensa \textit{(B.O.\ Nº 600, h.\ 3)}& \textbf{5.359 ha 77 a 83 ca}, según deslinde de Pfister y Martin aprobado por el Departamento Topográfico; base \$33.333,33, dos tercios de la tasación fiscal \\
\addlinespace
1916 & \textbf{Santa Mónica} y \textbf{Quitilipe} & --- & Deslinde, mensura y amojonamiento (\textbf{el segundo}, cuatro años después del de 1912) & Pedido por el doctor Carlos Serrey con poder de Manuel Goica; perito Herman Pfister\index{Pfister, Hermann}. Linderos: las cumbreras que colindan con los Yacones y San Alejo; al sud, «\textit{Getsmaní}» de Daniel Linares\index{Linares, Daniel}, separados por \textbf{una línea trazada por el ingeniero Pedro Cornejo el 9 de abril de 1889}, y propiedad de Abraham Estrada; al naciente el camino nacional a Jujuy; al norte el río de la Caldera; Quitilipe linda con F.\ Luis Acosta y María Arrieta \textit{(B.O.\ Nº 611, h.\ 2)}& --- \\
\addlinespace
1917 & \textbf{Potrero Gallinato} & Potrero de Gallinato & \textbf{Deslinde, mensura y amojonamiento} y \textbf{división de condominio} & Pedido por Camilo Cáceres ---impreso «\textit{Camilo Cacere}»---, \textbf{miembro de la comisión municipal del departamento desde quince días antes}, «\textit{con títulos bastantes}»; la finca «\textit{fue antes de Inocencio Aramayo}». Juez de feria Augusto F.\ Torino; perito Herman Pfister\index{Pfister, Hermann}. Linderos: norte, los filos de Santa Gertrudis, de los herederos de Ángel Solá; sud, los filos del cerro Mojotoro, que la divide de la propiedad de los Arias; este, \textbf{José Nogales}; oeste, \textbf{los herederos Yugra} \textit{(B.O.\ Nº 636, h.\ 4, y Nº 637, h.\ 4)}& --- \\
\addlinespace
1917 & Terreno en \textbf{La Calderilla} & La Calderilla & \textbf{Deslinde parcial} por el rumbo sud & De Juana Flores de Chaile; juez Augusto F.\ Torino; perito «\textit{Herman Pjister}» \textit{[sic]}. Linderos: norte, Juana Antonia Cáceres; sud, \textbf{Ildefonsa Burgos de Yugra}; este, \textbf{la familia Yugra}; oeste, el río de la Calderilla \textit{(B.O.\ Nº 664, h.\ 4 y 5)}& --- \\
\addlinespace
1917 & Terrenos en el partido de \textbf{Calderilla} & Calderilla & Deslinde, mensura y amojonamiento & Pedido por Augusto P.\ Matienzo con poder de \textbf{Abraham Fernández}, por sí y como tutor del menor Bernabé Fernández; juez David E.\ Gudiño; perito Herman Pfister\index{Pfister, Hermann}. Linderos: norte, \textbf{Teodoro Reyes}; sud, Eduardo Toconas; naciente, el cerro «\textit{Pucheta}», que los separa de la finca «\textit{Santa Gertrudis}»; poniente, el río de La Caldera \textit{(B.O.\ Nº 683, h.\ 4)}& --- \\
\addlinespace
1921 & \textbf{Cerro de Buena Vista} y \textbf{Sausal o Curuzú} & Partido de Vaqueros & Deslinde de las dos fincas & \textbf{Francisco S.\ Urquiza}\index{Urquiza, Francisco S.}, «de propiedad» suya; juez Daniel Etcheverry, agrim.\ Héctor Chiostri, secr.\ Tomás Izarrualde & Títulos de fs.\ 33 y 40; la resolución se imprime «\textit{Octubre 18 de 1020}» \textit{[sic]}; publicación ordenada «\textit{por una sola vez}» en el B.O.\ y treinta días en dos diarios. B.O.\ Nº 843, h.\ 8--9 \\
\addlinespace
1921 & \textbf{El Angosto} & Depto.\ La Caldera & \textbf{Deslinde, mensura y amojonamiento} & De \textbf{Antonia P.\ de Peralta}, por el procurador Elías Gallardo; juez Alberto Mendioroz, agrim.\ Rodolfo Chávez, secr.\ Tomás Izarrualde & \textbf{Linderos:} norte, herederos de «\textit{Luisa A.\ Costas}», hoy \textbf{Daniel Linares}\index{Linares, Daniel} y \textbf{Julio Mercado}; sud, el monte del \textbf{Quitilipe} separado por el río de La Caldera y propiedad de \textbf{Susana Aguilera}; este, \textbf{Augusto Reyes}; oeste, herederos de «\textit{Luis A.\ Costas}» \textit{[sic]}, hoy Daniel Linares. Aviso Nº 286. B.O.\ Nº 869, h.\ 5--6 \\
\addlinespace
1922 & Potrero del Gallinato & Potrero de Gallinato & Remate & --- & Base \$800; ejecución Alsina c/ José Nogales, martillero Leguizamón, fijado para el 28 de julio. \textit{No consta el resultado.} B.O.\ Nº 910, h.\ 16, aviso Nº 264 \\
\addlinespace
1923 & Finca «Calderilla» & La Calderilla & Remate & Mart.\ Leguizamón; juez Cánepa; juicio sucesorio de Silvano Murúa\index{Murúa, Silvano Ignacio}; con la novena parte de una casa de la capital y 217 vacunos; subasta fijada para el 11 de diciembre. \textit{No consta el resultado.} B.O.\ Nº 982, h.\ 11, aviso Nº 328 & Base \$7.000 \\
1924 & Finca «Calderilla» & La Calderilla & Remate de \textbf{derechos y acciones} & Mart.\ Leguizamón; juez Figueroa; Marcos Alsina c/ sucesión de los esposos Mayo (incidente de honorarios); derechos de Rafaela de Mamani, Pedro Ríos y Manuel Lozano; subasta el 16 de abril. B.O.\ Nº 1001, h.\ 18 & Base \$533,33 \\
1924 & Finca «Calderilla» & La Calderilla & Remate de \textbf{la mitad} & \textit{Misma ejecución y mismo martillero, tres meses después, con otros titulares ---Pedro Ruiz y Manuel Lozano--- y otro objeto. \textbf{No se concilia con la fila anterior.}} Subasta el 11 de agosto. B.O.\ Nº 1016, h.\ 15, aviso 633 & Base \$266,66 \\
\addlinespace
1924 & \textbf{Vaqueros o Entre Ríos} & Partido de Vaqueros & Deslinde (auto del 5 de diciembre de 1923) & Natal López Cross; lindero este, Francisco Urquiza\index{Urquiza, Francisco S.}; oeste, el camino nacional; agrim.\ Pedro J.\ Frías & --- \\
\addlinespace
1924 & Finca de pastoreo & Quebrada de Mojotoro & Remate \textbf{dos veces} & Mart.\ Decavi; lindero este, Silvano Murúa en el primer aviso y Desiderio Gumiel en el segundo; norte, la cumbre que la separa del potrero de Gallinato & 2.200 $\times$ 3.500 m; base \$1.500, dos tercios de la avaluación \\
\addlinespace
1924 & \textbf{Wierna} & --- & Deslinde & Silvano Murúa y María Fanny Ovejero de Murúa; lindero norte, la finca «Caldera» de Daniel Linares\index{Linares, Daniel}; este, el río de La Calderilla & --- \\
\addlinespace
1925 & El Sauce o Paraíso, y fracción Ceival & Partido El Sauce & Deslinde, mensura y amojonamiento & Linderos: los Murúa, los Patrón, los «Noques» de Linares & --- \\
\addlinespace
1926 & Los Peñones & Peñones & \textbf{Reposición de título} & María Isabel Gómez; información sumaria ante el Juez de Paz & --- \\
\addlinespace
1926 & Finca rural sin nombre & Margen del río Wierna & Deslinde & Sándalio Viveros; linderos Toranzos, Arquati & --- \\
\addlinespace
1945 & Finca «Las Lagunas» & Las Lagunas & \textbf{Posesión treintañal} & Edicto repetido en más de veinte ediciones. Límites de naciente a poniente: «\textit{desde las cumbres altas de las serranías de la Caldera hasta dar con el río de los Yacones, enfrentando con el arroyo de las carreras}». Se oficia al juez de paz de La Caldera y se da intervención al Fiscal de Gobierno. \textit{No consta el resultado.} B.O.\ Nº 2202 h.\ 6 y siguientes & --- \\
1945 & Terreno en La Calderilla & La Calderilla & \textbf{Posesión treintañal} & Frente de \textbf{108 metros con 20 cm}; linda con propiedad de la \textbf{sucesión de Micaela de Valencia} y, al oeste, con el \textbf{río de La Caldera}. Edicto repetido en más de veinte ediciones. \textit{No consta el resultado.} B.O.\ Nº 2403 h.\ 7 y siguientes & --- \\
1927 & \textbf{Los Sauces} & Los Sauces & Remate \textbf{tres veces} & Mart.\ Roberto del Carril; juez Ángel María Figueroa; \textbf{juicio sucesorio de Domingo Royo}, miembro de la Comisión Municipal en 1922--1923. Base \$50.000 el 3 de junio, \$14.000 el 8 de julio y el mismo aviso el 26 de agosto, con tres números de aviso distintos. \textit{No consta el resultado.} B.O.\ Nº 1169 h.\ 14, Nº 1174 h.\ 11 y Nº 1182 h.\ 14 & \textbf{Base \$50.000}, bajada a \$14.000 \\
1927 & Puesto de las Garrapatas & --- & Deslinde & Silvano I. Murúa; linda con Los Porongos de Daniel Linares\index{Linares, Daniel} & --- \\
\addlinespace
1927 & Los Sauces & Hasta el límite con Jujuy & Remate (sucesión Domingo Royo) & Mart.\ Forcada; linda al Sud con Angostura de Daniel Linares & 8.000 $\times$ 17.000 m; base \$50.000 y luego \$14.000 \\
\addlinespace
1927 & Casa quinta & Lesser & Remate (Banco de la Nación) & Mart.\ Forcada; tres fracciones unidas & 10 ha; base \$6.000 \\
\addlinespace
1929 & Angostura o Angosto de Arias & --- & Deslinde & Linda con Los Sauces, con Santa Rufina de Bernardo Austerlitz y con Julio Forcada & --- \\
\addlinespace
1929 & Casa quinta & Vaqueros & Remate & Mart.\ Peñalba Herrera; linda al Norte con la señora de Bianchi y al Sud con el río Vaqueros & 5 ha; base \$1.000, dos tercios de la avaluación \\
\addlinespace
1930 & Estancia \textbf{Potrero de Gallinato} & Potrero de Gallinato & Remate (Rocha c/ González) & Mart.\ Leguizamón; se remata junto con una casa de la capital & Base \$1.333,33 \textit{(cifra dañada en el original)} \\
\addlinespace
1930 & Finca y ganado & --- & Remate & Mart.\ Forcada; lindero Norte, herederos de Julio Mercado & Base \$900; sin superficie ni catastro \\
\addlinespace
1932 & \textbf{Las Garrapatas, Cañada Seca o Peñones, El Durazno o Peñones, Los Matos y Severino} & --- & Deslinde de \textbf{cinco fincas} a la vez & Silvano Y.\ Murúa y Fanny Ovejero Lacroix de Murúa & --- \\
\addlinespace
1932 & Casa en el pueblo & Pueblo & Remate \textbf{tres veces} & Mart.\ Salvatierra; Junco de Bertini c/ Cornejo & \$3.333,33 $\rightarrow$ \$2.500 $\rightarrow$ \textbf{sin base} \\
\addlinespace
1933 & \textbf{La Despensa} & Pueblo de La Caldera & Remate \textbf{voluntario} & Mart.\ Rossi; por orden de Alberto, Abel y Adolfo Murúa; linderos, la sucesión y los herederos de Silvano Murúa & \textbf{2.700 ha}; base \$35.000, avaluación \$50.000 \\
\addlinespace
1934 & Puesto Despensa & --- & Deslinde & Alberto, Adolfo y Abel Murúa; perito Hermann Pfister; lindero Norte, Francisco Sosaya en parte y el puesto Cañada Seca o Peñones de Silvano Murúa & \textbf{$\approx$ 2.600 ha} \\
\addlinespace
1934 & El Angosto & Partido del Angosto & Remate \textbf{tres veces} & Mart.\ Figueroa Echazú & \$10.000 $\rightarrow$ \$6.500 $\rightarrow$ \textbf{sin base} \\
\addlinespace
1934 & San Jorge y San Félix & --- & Remate \textbf{dos veces} & Mart.\ Leguizamón & \$8.000 + \$8.000 por separado $\rightarrow$ \textbf{\$6.000 unidas} \\
\addlinespace
1935 & Vaqueros o Entre Ríos & Vaqueros & Remate hipotecario del Banco Provincial & c/ Ramón Bracheri; lindero este, Francisco Urquiza\index{Urquiza, Francisco S.} & Base \$12.000 \\
\addlinespace
1936 & \textbf{La Despensa} & Partido del Sauce \textit{(el aviso de 1933 dice pueblo de La Caldera)} & Remate \textbf{dos veces} & Mart.\ Esquiú; Salvador Rosa c/ Abel, Adolfo y Alberto Murúa & Base \$23.333,32 y luego \textbf{\$25.000} \\
\addlinespace
1937 & \textbf{Chalchanio} & --- & Remate & Mart.\ Decavi; Terán c/ Ramón I. Juárez\index{Juárez, Ramón} & Base \$2.333,33. \textbf{Folio 342, asiento 438, Libro B} \\
\addlinespace
1937 & \textbf{La Helvecia}, antes Calderilla & La Calderilla & Deslinde & Hermann Pfister\index{Pfister, Hermann}; \textbf{lindero Norte, Augusto Regis}\index{Regis, Augusto} & --- \\
\addlinespace
1938 & Abra de Leser \textit{(así en el edicto)} & La Calderilla & Deslinde & Dr.\ Lucio Ortiz; linda al Sud con la finca \textbf{Lesser, de Luis Patrón Costas} & --- \\
\addlinespace
1939 & Terreno del pueblo & Pueblo & \textbf{Posesión treintañal} & Curia Eclesiástica de Salta; \textbf{linderos Norte y Oeste, Dr.\ Carlos Serrey}\index{Serrey, Carlos} & 126,80 m de frente \\
\addlinespace
1942 & Casa y terreno, quinta parte indivisa & La Calderilla & Remate & Mart.\ Leguizamón; sucesorio de Susana Aguilera de Ruiloba & Base \$500 \\
\addlinespace
1942 & \textbf{Yacones o Tránsito} & Partido de La Calderilla & Deslinde & Dr.\ Lucio Ortiz; límites, el arroyo de los Yacones, los ríos Yacones y Wierna y la serranía alta; agrim.\ José F.\ Campilongo & --- \\
\addlinespace
1943 & Predio de la Comisaría & Pueblo & \textbf{Posesión treintañal} & \textbf{El Gobierno de la Provincia}, que ocupaba sin título desde 1913 & 8.000 m\textsuperscript{2} \\
\addlinespace
1945 & Las Lagunas & Las Lagunas & \textbf{Posesión treintañal} & Gregorio Maidana; título Bejarano a Arroyo, en el que Maidana no figura & \textbf{Sin superficie ni catastro} \\
\addlinespace
1945 & \textbf{Chalchanio}, cuarta parte indivisa & --- & Remate & Mart.\ Rivas; Strachan, Yáñez y Cía.\ c/ Martín Herrera & 490 ha 9 a 36 cm\textsuperscript{2}; base \$671,01. \textbf{Mismo folio 342} \\
\addlinespace
1945 & Terreno con lo plantado & La Calderilla & \textbf{Posesión treintañal} & Hermann Pfister\index{Pfister, Hermann}; lindero Oeste, el río de La Caldera & 108,28 m de frente por 2.395 m de fondo \\
\bottomrule
\end{longtable}}

\section*{Lo que la tabla muestra}

\textbf{Primero: la partición de la sucesión Murúa es el hecho dominial mayor del período.} En
1924 Silvano Ignacio Murúa\index{Murúa, Silvano Ignacio} y su esposa, María Fanny Ovejero, hacen deslindar la finca
Wierna; en 1927 el mismo Silvano ---«Silvano I.» en el edicto--- hace deslindar el Puesto de las
Garrapatas; en 1932, él ---«Silvano Y.» esta vez--- y su esposa hacen deslindar \textbf{cinco fincas de una vez} ---Las Garrapatas, Cañada Seca o
Peñones, El Durazno o Peñones, Los Matos y Severino---. En octubre de 1933 Alberto, Abel y
Adolfo Murúa ponen en remate \textbf{por su cuenta} La Despensa, dos mil setecientas
hectáreas, con base de \$35.000 sobre una avaluación de \$50.000, y el aviso nombra ya como
vecinos a la sucesión y a los herederos del padre de Silvano Ignacio, también Silvano, muerto antes de agosto de 1927. En 1934 se deslinda el Puesto
Despensa, de unas dos mil seiscientas hectáreas; y en 1936 La Despensa se remata en la
ejecución de Salvador Rosa contra esos mismos tres herederos, con los otros herederos como
linderos. Deslinde, intento de venta, partición, ejecución: el ciclo completo, en doce años y
sobre el mismo patrimonio. Y una de
esas cinco fincas de 1932 se llama \textbf{El Durazno}.

\textbf{Segundo: los remates caen.} Ocho bienes salen a subasta más de una vez en un mismo
juicio ---la sucesión Catacata, la finca de pastoreo de Mojotoro, Los Sauces, la casa del
pueblo, El Angosto, San Jorge y San Félix, La Despensa y \textbf{San Alejo y Santa Rufina}---,
y en cinco de ellos la base baja. \textbf{El caso extremo es San Alejo}: cinco llamados entre
diciembre de 1911 y octubre de 1914, con la base por hectárea cayendo de \$10,07 a \$3,20
---un sesenta y ocho por ciento--- sin que el archivo registre venta (capítulo
\ref{cap:siglo}). En los otros cuatro: Los Sauces en 1927 (de \$50.000 a
\$14.000), la casa del pueblo en 1932 (\$3.333,33, \$2.500, sin base), El Angosto en 1934
(\$10.000, \$6.500, sin base) y San Jorge con San Félix el mismo año (de \$16.000 por separado
a \$6.000 unidas). En dos de ellos, la base desaparece. \textbf{En 1932 y 1934 no había quien comprara
tierra en este departamento al precio de tasación.} La única base que sube es la de La Despensa
en 1936, de \$23.333,32 a \$25.000 ---cifra que el capítulo \ref{cap:fincas} pide cotejar contra el original antes de citarla---, que es lo contrario de lo esperable en un remate reiterado y
que este libro no explica. Un dato lo enmarca sin explicarlo: aquella primera base de 1936 es,
salvo un centavo, dos tercios de los \$35.000 que los propios Murúa habían pedido en 1933.

\textbf{Tercero: el Estado también tuvo que ir al juez.} Cuatro de las cincuenta y ocho operaciones de esta tabla son
juicios de posesión treintañal, y uno de ellos lo promueve \textbf{el Gobierno de la
Provincia}, en 1943, sobre los ocho mil metros cuadrados donde funcionaba su propia Comisaría
desde 1913 o antes. Los otros tres los promueven la Curia Eclesiástica (1939), Gregorio
Maidana (1945) y Hermann Pfister\index{Pfister, Hermann} (1945). \textbf{La vía para regularizar tierra en este
departamento era la misma para el Estado, para la Iglesia y para los particulares}: treinta años
de ocupación y un edicto.

\textbf{Cuarto: describir por linderos no fue una etapa que se superó.} El primer aviso del
corpus que cita folio, asiento y libro es el de Chalchanio en 1937. Ocho años después, el
edicto de Las Lagunas identifica una finca entera por cursos de agua y vecinos, sin superficie
ni catastro. \textbf{Las dos formas conviven hasta el final del período}, y conviven con la
tercera que el capítulo \ref{cap:ribera} encuentra setenta años después.

\textbf{Y quinto: los nombres se repiten, y son los mismos del resto del libro.} Daniel
Linares\index{Linares, Daniel} aparece como lindero en 1924, 1925, 1927 y 1927. Bernardo
Austerlitz\index{Austerlitz, Bernardo}, en 1929. Francisco S.\
Urquiza\index{Urquiza, Francisco S.}, en 1924 y en 1935. Silvano Murúa, en 1924. Augusto Regis\index{Regis, Augusto}, en 1937.
Carlos Serrey\index{Serrey, Carlos}, en 1939. Luis Patrón Costas, en 1938. Los herederos de
Julio Mercado, en 1930 ---un Julio Mercado integraba en 1928 la Comisión Municipal (capítulo \ref{cap:politica}), y el apellido Mercado vuelve en 1972 entre los expropiados para el
embalse. \textbf{No hace falta inferir nada: los linderos de las fincas del departamento son,
uno por uno, los nombres que el apéndice \ref{cap:personas} sigue por otras vías.}

\section*{Cómo se usa}

Cada fila da tres cosas que un pedido necesita: \textbf{el nombre con que la finca se llamaba
entonces}, \textbf{el año del acto} y \textbf{el juzgado o el martillero}. Con eso se localiza
el expediente; con el expediente, el folio; con el folio, la cadena. Para dos de las
cincuenta y ocho filas el propio aviso ya trae el folio, y son la misma: \textbf{Chalchanio, folio
342, asiento 438 del Libro B}, citado en 1937 y otra vez en 1945 con distinto ejecutado.

Falta, y conviene decirlo, \textbf{la correspondencia entre estos nombres y las matrículas
actuales}. El capítulo \ref{cap:fincas} la pide desde los trece partidos de 1909; con esta tabla
y con el padrón de parajes de 1940 que ese mismo capítulo transcribe, el pedido puede formularse
sobre una lista cerrada de nombres en lugar de sobre una categoría.
